\documentclass[11pt,a4paper]{article}
\usepackage[margin=2.5cm]{geometry}
\usepackage[T1]{fontenc}
\usepackage[utf8]{inputenc}
\usepackage[spanish,es-tabla]{babel}
\usepackage{lmodern,microtype,booktabs,xcolor,hyperref,amsmath,enumitem}
\hypersetup{colorlinks=true,linkcolor=blue!50!black,urlcolor=blue!50!black}
\setlength{\emergencystretch}{4em}
\newcommand{\acc}{\mathrm{acc}}

\title{Validación de los clusters por $\Delta$ y de la búsqueda de arquitecturas:\\ PCA, transferencia, FEMNIST, DBSCAN, interacción arquitectura $\times$ nodo y P3}
\author{Proyecto EvoFederated}
\date{13 de noviembre de 2026}

\begin{document}
\maketitle

\begin{abstract}
\noindent Se valida la agrupación de nodos federados por similitud de las actualizaciones $\Delta$ de un modelo sonda (k-means esférico, $C=10$, $N=100$ nodos) y, a partir de ella, la búsqueda de arquitecturas, con seis estudios que no repiten el NAS completo salvo el último. \textbf{(1)} Aplicar PCA antes del clustering es igual o peor que usar $\Delta$ crudo; la pérdida viene de centrar, no de reducir dimensiones. \textbf{(2)} En CIFAR-10 con Dirichlet, un modelo entrenado en un nodo rinde entre 10 y 35 puntos de accuracy más en nodos de su cluster que en otros (10 de 10 semillas significativas en los tres $\alpha$); $\Delta$ recupera entre el 76\,\% y el 89\,\% de la ganancia de una agrupación con etiquetas. \textbf{(3)} En FEMNIST, con partición natural por escritor, el efecto es pequeño pero real (3.3 puntos, 20 de 20 semillas) y $\Delta$ iguala a la agrupación con etiquetas; de forma exploratoria, $\Delta$ conserva una señal de $\approx$1 punto al controlar por la distribución de etiquetas. \textbf{(4)} Agrupar por \texttt{hsf\_id} transfiere casi tan bien como $\Delta$ pero son agrupaciones distintas, y DBSCAN, con o sin PCA, no mejora a k-means: encuentra 2--6 clusters grandes con 17--43\,\% de ruido. \textbf{(5)} Con un modelo global FedAvg por arquitectura, existe una interacción arquitectura $\times$ nodo estable pero pequeña (SD de 4.7 puntos frente a 18.1 del efecto de arquitectura); el ranking de arquitecturas es casi el mismo en todos los nodos (Spearman 0.92 entre nodos), elegir por nodo o por cluster no supera a la mejor global y bastan pocos nodos para elegir (10 nodos recuperan el 98.5\,\% de la mejora posible). \textbf{(6)} La silueta elige $k=2$ en las 20 federaciones; con ese $k$, NSGA-III con un objetivo por cluster y accuracy penalizada por tamaño no se distingue de NSGA-II con $C=10$ ni de la búsqueda aleatoria (hipervolumen 0.515, 0.518 y 0.499; todos los intervalos incluyen 0, 3 semillas) y usa 4.6 veces menos entrenamientos locales; la evolución mejora a la población pero no al mejor individuo. Los puntos (3) exploratorio, (4) y los análisis de la Sección~\ref{sec:p3} sobre evolución por generación son post hoc; (5) y (6) se preregistraron.
\end{abstract}

\section{Motivación y medidas}
Hasta ahora la calidad de los clusters se medía con $\eta^2$ y ARI, que dependen de la distribución de etiquetas y son un sustituto de lo que importa: que un representante hable por su cluster. Aquí se usan dos medidas.
\begin{itemize}[leftmargin=*,itemsep=2pt]
\item \textbf{$\eta^2$ y ARI.} $\eta^2$ es la fracción de la varianza de la distribución de clases entre nodos explicada por los clusters (0 a 1). El ARI mide el acuerdo con el oráculo, que agrupa con k-means sobre la distribución de clases real (información no privada, solo referencia).
\item \textbf{Transferencia.} $A[s,j]$ es la accuracy en el conjunto de validación del nodo $j$ de un modelo entrenado solo en el nodo $s$ (arquitectura distinta de la sonda, misma inicialización para todas las fuentes). La ganancia es
\[
\mathrm{gap}=\frac1N\sum_s\Bigl(\overline{A[s,j\in c(s),j\neq s]}-\overline{A[s,j\notin c(s)]}\Bigr).
\]
Por semilla se obtiene un $p$ por permutación de las etiquetas de cluster (10\,000, tamaños conservados); entre semillas se usa Wilcoxon pareado con corrección de Holm e intervalos bootstrap.
\end{itemize}
Las hipótesis, el estadístico y los criterios de decisión de los estudios de transferencia se escribieron antes de correr (\texttt{docs/TRANSFER\_DESIGN.md}). Los análisis marcados como exploratorios no estaban en ese preregistro.

\section{PCA frente a $\Delta$ crudo}
CIFAR-10, 30 semillas por celda, solo calidad de clusters. La Tabla~\ref{tab:pca} resume $\eta^2$ y ARI con el oráculo.

\begin{table}[h]
\centering\small
\begin{tabular}{llccccc}
\toprule
$\alpha$ & métrica & crudo & centrado & PCA 95\,\% & PCA 80\,\% & PCA $k{=}5$\\
\midrule
0.1 & $\eta^2$ & \textbf{0.696} & 0.510 & 0.520 & 0.518 & 0.395\\
0.1 & ARI oráculo & \textbf{0.609} & 0.307 & 0.316 & 0.320 & 0.229\\
0.5 & $\eta^2$ & 0.468 & 0.447 & 0.450 & 0.452 & 0.410\\
0.5 & ARI oráculo & 0.311 & 0.287 & 0.292 & 0.289 & 0.231\\
1.0 & $\eta^2$ & 0.424 & 0.404 & 0.413 & 0.418 & 0.381\\
1.0 & ARI oráculo & 0.252 & 0.232 & 0.240 & 0.248 & 0.212\\
\bottomrule
\end{tabular}
\caption{PCA antes del clustering (media de 30 semillas). ``Centrado'' resta la media sin reducir dimensiones.}
\label{tab:pca}
\end{table}

Con heterogeneidad fuerte ($\alpha=0.1$) el crudo gana por mucho (brecha de 0.19 en $\eta^2$); con $\alpha=0.5$ y $1.0$ la diferencia es de 0.01 a 0.03, sin prueba pareada. ``Centrado sin reducir'' rinde igual que el PCA, así que la pérdida viene del centrado; reducir demasiado ($k=5$) empeora siempre. Por qué el centrado perjudica no está comprobado.

\section{Transferencia entre nodos: CIFAR-10 con Dirichlet}
Diez semillas por $\alpha$, cinco épocas por fuente, $N=100$, $C=10$. El gap de una agrupación aleatoria es cero por construcción (Tabla~\ref{tab:tr}).

\begin{table}[h]
\centering\small
\begin{tabular}{cccccc}
\toprule
$\alpha$ & acc.\ media & gap $\Delta$ (IC95\,\%) & gap oráculo & $\Delta$/oráculo & Holm-$p$\\
\midrule
0.1 & 0.134 & \textbf{0.348} (0.335--0.362) & 0.394 & 0.89 & 0.003\\
0.5 & 0.170 & \textbf{0.166} (0.152--0.179) & 0.210 & 0.81 & 0.003\\
1.0 & 0.182 & \textbf{0.100} (0.092--0.106) & 0.133 & 0.76 & 0.003\\
\bottomrule
\end{tabular}
\caption{Transferencia en CIFAR-10. En los tres $\alpha$ las 10 semillas tienen $p<0.05$ por permutación; $\Delta$/oráculo es la mediana entre semillas.}
\label{tab:tr}
\end{table}

La hipótesis H1 se apoya en los tres $\alpha$. Holm-$p=0.003$ es el mínimo alcanzable con 10 semillas (todas en la misma dirección). El gap mide sobre todo el solapamiento de clases entre nodos: con Dirichlet, la similitud funcional depende casi solo de las etiquetas. De forma exploratoria, el coeficiente de ``mismo cluster'' en la regresión $A\sim\text{mismo cluster}+d_{\text{Hellinger}}$ es 0.077, 0.075 y 0.048.

\section{Réplica en FEMNIST (partición natural)}
Cada cliente es un escritor, sin Dirichlet: 100 escritores por semilla (con al menos 100 muestras), 20 semillas, mediana de 38 de 62 clases por nodo (clase dominante 9\,\%) y 92 muestras de entrenamiento. Con 5 épocas los modelos quedaban casi en azar, así que se fijaron \textbf{50 épocas} con una semilla fuera del estudio, según la accuracy en el propio nodo (5 épocas: 6.6\,\%; 20: 7.9\,\%; 50: 17.3\,\%; azar: 1.6\,\%) y no según el gap.

\begin{table}[h]
\centering\small
\begin{tabular}{ccccc}
\toprule
gap $\Delta$ (IC95\,\%) & gap oráculo & gap aleatorio & semillas $p<0.05$ & Wilcoxon $\Delta$ vs azar\\
\midrule
\textbf{0.033} (0.030--0.035) & 0.030 & 0.0002 & 20/20 & $9.5\times10^{-7}$\\
\bottomrule
\end{tabular}
\caption{Transferencia en FEMNIST (accuracy media 0.097).}
\end{table}

$\Delta$ supera al oráculo en 16 de 20 semillas (diferencia media $+0.0026$). El efecto es pequeño porque cada nodo tiene pocas muestras y 62 clases. \textbf{Exploratorio:} al controlar por la distancia de Hellinger entre etiquetas, el coeficiente de ``mismo cluster'' es 0.0101 (IC95 bootstrap 0.0046--0.0152; positivo en 16/20 semillas; Wilcoxon $p=0.0018$), es decir, $\Delta$ agrupa nodos funcionalmente parecidos más de lo que explican sus etiquetas. Es un modelo lineal sobre una relación que no lo es, por lo que puede quedar confusión residual.

\section{\texttt{hsf\_id} frente a $\Delta$ en FEMNIST (exploratorio)}
El \texttt{hsf\_id} es la serie de formularios NIST; cada escritor pertenece a una de 7. Se reutilizan las mismas matrices $A$ (Tabla~\ref{tab:hsf}).

\begin{table}[h]
\centering\small
\begin{tabular}{lcc}
\toprule
agrupación & gap (IC95\,\%) & semillas $p<0.05$\\
\midrule
\texttt{hsf\_id} (7 grupos) & 0.0307 (0.0288--0.0325) & 20/20\\
$\Delta$, $C=10$ & 0.0327 (0.0301--0.0354) & 20/20\\
$\Delta$, $C=7$ & \textbf{0.0348} (0.0319--0.0375) & 20/20\\
\bottomrule
\end{tabular}
\caption{Agrupar por \texttt{hsf\_id} frente a $\Delta$.}
\label{tab:hsf}
\end{table}

Con $C=7$, $\Delta$ supera a hsf en 16 de 20 semillas (Wilcoxon bilateral $p=0.0002$); con $C=10$, en 11 de 20 ($p=0.076$). El ARI entre $\Delta$ y hsf es 0.14 ($C=10$) y 0.17 ($C=7$): son agrupaciones distintas. En la regresión conjunta $A\sim\text{mismo}_\Delta+\text{mismo}_{\text{hsf}}+d_{\text{Hellinger}}$, los coeficientes son 0.0100 (0.0045--0.0150) para $\Delta$ y 0.0059 (0.0051--0.0068) para hsf. Los tamaños de cluster desiguales de $\Delta$ frente a los grupos hsf, parecidos entre sí, no se controlaron.

\section{DBSCAN, con y sin PCA (exploratorio)}
Mismas matrices $A$, sin reentrenar (CIFAR-10: 3 $\alpha\times$10 semillas; FEMNIST: 20 semillas). Se usan $\Delta$ crudo, PCA 95\,\% y PCA 80\,\%; DBSCAN con distancia coseno, \texttt{min\_samples}$=5$ y $\varepsilon$ igual al cuantil $q\in\{0.3,0.5,0.7,0.9\}$ de la distancia al quinto vecino. La variante principal es la de mayor silueta (criterio no supervisado); el ruido se asigna al centroide más cercano. La comparación justa es con k-means del \emph{mismo} $k$ y misma representación, porque el gap depende de $k$. La reconstrucción de $\Delta$ reproduce los clusters guardados (ARI $=1$).

\begin{table}[h]
\centering\small
\begin{tabular}{llcrrrrrr}
\toprule
datos & $\alpha$ & repr. & válidas & $k$ & ruido & gap DBSCAN & gap k-means & Holm-$p$\\
\midrule
CIFAR-10 & 0.1 & crudo & 10 & 5.7 & 25\,\% & 0.217 & 0.237 & 0.54\\
CIFAR-10 & 0.1 & PCA 95 & 10 & 2.5 & 40\,\% & 0.054 & 0.033 & 0.041\\
CIFAR-10 & 0.1 & PCA 80 & 10 & 3.4 & 38\,\% & 0.088 & 0.069 & 0.26\\
CIFAR-10 & 0.5 & crudo & 5 & 2.6 & 25\,\% & 0.054 & 0.072 & --\\
CIFAR-10 & 0.5 & PCA 95 & 7 & 3.0 & 17\,\% & 0.050 & 0.081 & 0.063\\
CIFAR-10 & 0.5 & PCA 80 & 9 & 2.8 & 20\,\% & 0.049 & 0.076 & 0.016\\
CIFAR-10 & 1.0 & crudo & 4 & 2.8 & 23\,\% & 0.031 & 0.047 & --\\
CIFAR-10 & 1.0 & PCA 95 & 9 & 2.3 & 18\,\% & 0.024 & 0.043 & 0.016\\
CIFAR-10 & 1.0 & PCA 80 & 8 & 2.6 & 17\,\% & 0.025 & 0.042 & 0.023\\
FEMNIST & natural & crudo & 20 & 2.05 & 34\,\% & 0.0555 & 0.0551 & 0.54\\
FEMNIST & natural & PCA 95 & 20 & 2.05 & 30\,\% & 0.0553 & 0.0553 & 0.97\\
FEMNIST & natural & PCA 80 & 20 & 2.0 & 29\,\% & 0.0552 & 0.0548 & 0.65\\
\bottomrule
\end{tabular}
\caption{DBSCAN (variante de mayor silueta) frente a k-means con el mismo $k$. ``Válidas'': semillas con al menos una variante no degenerada (de 10 en CIFAR-10 y 20 en FEMNIST). Holm sobre las 4 celdas de cada representación; ``--'': muy pocas semillas válidas para el test.}
\label{tab:db}
\end{table}

DBSCAN encuentra pocos clusters (2.3 a 5.7 en CIFAR-10, 2.0 en FEMNIST) con 17--43\,\% de ruido, y con $q$ alto todo se fusiona en un solo cluster: entre el 48\,\% y el 90\,\% de las corridas (semilla $\times\,q$) de CIFAR-10 son degeneradas. Frente a k-means del mismo $k$, DBSCAN empata (CIFAR-10 con $\alpha=0.1$ crudo; FEMNIST) o pierde (CIFAR-10 con $\alpha=0.5$ y $1.0$ con PCA, Holm-$p$ 0.016--0.023); la única celda a su favor es CIFAR-10 con $\alpha=0.1$ y PCA 95\,\% ($p=0.041$), con gaps muy bajos. Frente a k-means con $C=10$ los contrastes salen significativos pero no son interpretables, porque $k$ difiere (en CIFAR-10 DBSCAN pierde; en FEMNIST, con $k=2$, el gap es 0.055 frente a 0.033). Para un NAS con unos 10 representantes, DBSCAN da 2--6 clusters grandes, lo que implica mayor varianza del estimador; no se integra como método de muestreo.

\section{Interacción arquitectura $\times$ nodo (FEMNIST, modelo global FedAvg)}
\label{sec:an}
La pregunta es si hay arquitecturas que funcionan mejor en unos nodos que en otros o si rinden igual en todos. Importa para el NAS: sin interacción, el ranking global basta y los objetivos por cluster o nodo no aportan. Se entrenan 24 arquitecturas aleatorias del espacio v2 (hasta 3 bloques convolucionales cuyos slots pueden apagarse, hasta 2 capas totalmente conectadas, pooling $2\times2$) con FedAvg sobre los 100 nodos (20 rondas $\times$ 2 épocas, Adam, $\mathrm{lr}=10^{-3}$, elegido con un piloto en una semilla fuera del estudio: 28.2\,\% frente a 22.1\,\% de $10^{-2}$), en 3 federaciones y con 2 réplicas de entrenamiento (otra inicialización y otro orden de datos). Se mide la accuracy de cada modelo en el val y en el test de cada nodo.

\textbf{Estadístico.} Con $Y[a,n]$ la accuracy de la arquitectura $a$ en el nodo $n$ y $R$ su doble centrado ($Y$ menos medias por arquitectura y por nodo, más la media global), la interacción estable es $V=\overline{R_{\mathrm{val}}^{(r)}\cdot R_{\mathrm{test}}^{(r')}}$ con $r\neq r'$. Es un estimador sin sesgo de la varianza de la interacción verdadera, porque el ruido de entrenamiento y el de evaluación son independientes entre los dos factores. Se prueba con 10\,000 permutaciones de los nodos y se acompaña de un intervalo bootstrap sobre nodos y de un control en escala logit.

\begin{table}[h]
\centering\small
\begin{tabular}{lcc}
\toprule
 & accuracy & logit\\
\midrule
SD del efecto de \textbf{arquitectura} & 18.1 pp & 1.15\\
SD del efecto de \textbf{nodo} & 7.4 pp & 0.47\\
SD de la \textbf{interacción estable} & \textbf{4.7 pp} & 0.17\\
fracción de la variación residual que se repite & 49\,\% & 20\,\%\\
IC95 de $V>0$ & 3/3 semillas & 3/3\\
$p$ por permutación & 0.0001 en cada semilla & 0.0001\\
\bottomrule
\end{tabular}
\caption{Interacción arquitectura $\times$ nodo (media de 3 federaciones; 0.0001 es el mínimo alcanzable con 10\,000 permutaciones).}
\label{tab:an}
\end{table}

La interacción existe y se repite con otro entrenamiento y otra muestra del mismo nodo, también en logit, pero es pequeña frente al efecto de la arquitectura (su varianza es $\approx$7\,\% de la de la arquitectura); parte del efecto en accuracy es de escala (49\,\% baja a 20\,\% en logit). El ranking de arquitecturas es casi el mismo en todos los nodos: Spearman 0.97 entre réplicas de entrenamiento, 0.99 entre federaciones y 0.92 entre el ranking de un nodo y el de otro.

\textbf{Consecuencia práctica.} Se elige la arquitectura con el val de una réplica y se evalúa con el test de la otra, comparando con la mejor global elegida con todos los nodos. Elegir \emph{por nodo} cambia la accuracy en $-1.25$ pp y elegir \emph{por cluster} $\Delta$ en $-0.52$ pp; en ningún caso (0 de 3 semillas) el IC95 queda por encima de 0. Con 36--54 muestras por nodo, la señal de la interacción es más débil que el ruido de seleccionar.

\begin{table}[h]
\centering\small
\begin{tabular}{lcc}
\toprule
nodos usados para elegir & de la mejora posible sobre el azar & pierde frente a la mejor\\
\midrule
1 & 91.2\,\% & $-2.8$ pp\\
3 & 96.5\,\% & $-1.1$ pp\\
10 & 98.5\,\% & $-0.5$ pp\\
30 & 99.9\,\% & $-0.04$ pp\\
\bottomrule
\end{tabular}
\caption{Elegir la arquitectura con $k$ nodos y evaluarla en el resto (val de una réplica, test de la otra).}
\label{tab:pick}
\end{table}

\textbf{¿Hace falta elegir arquitectura?} Sí: la accuracy media entre las 24 arquitecturas va de $\approx$5\,\% (la peor) a $\approx$25\,\% (la mediana) y $\approx$58\,\% (la mejor), y solo una queda a menos de 2 pp de la mejor. Pero basta evaluar con pocos nodos (Tabla~\ref{tab:pick}). Accuracy y parámetros tienen un Spearman de 0.51; 8 de las 24 arquitecturas están en el frente de Pareto (de 8\,062 parámetros y 5.4\,\% a 69\,214 parámetros y 57.9\,\%), y la mejor usa 69 mil de los 354\,494 del espacio.

\section{P3: $k$ óptimo y NSGA-III con un objetivo por cluster}
\label{sec:p3}
\textbf{Propuesta.} Elegir el número de clusters $k\le10$, usar un objetivo por cluster (la accuracy del nodo evaluador de ese cluster, con un representante re-sorteado en cada generación) y meter el tamaño dentro de cada objetivo, de modo que el número de objetivos es $k$: $\mathrm{obj}_c=\acc_c-\lambda\,p/p_{\mathrm{ref}}$, con $\lambda=1$, $p$ el número de parámetros y $p_{\mathrm{ref}}=354\,494$ el máximo del espacio. Es E5 (NSGA-III con un objetivo por cluster) más la elección de $k$ y el tamaño dentro del objetivo. El criterio de $k$ es la silueta coseno de k-means esférico sobre $\Delta$ (no supervisado), fijado antes de correr.

\textbf{Elección de $k$.} La silueta decrece de forma monótona con $k$, así que elige $k=2$ en las 20 federaciones (Tabla~\ref{tab:k}); el gap de transferencia también es mayor con $k=2$. En la práctica, P3 fue NSGA-III con \textbf{2 objetivos}, no con hasta 10: la idea de un objetivo por cluster con muchos clusters no se puso a prueba. Es coherente con DBSCAN, que también encontraba $\approx$2 clusters.

\begin{table}[h]
\centering\small
\begin{tabular}{lccccccc}
\toprule
$k$ & 2 & 3 & 4 & 5 & 6 & 8 & 10\\
\midrule
silueta & \textbf{0.291} & 0.234 & 0.118 & 0.089 & 0.071 & 0.059 & 0.039\\
gap de transferencia & 0.055 & 0.053 & 0.039 & 0.037 & 0.035 & 0.034 & 0.033\\
\bottomrule
\end{tabular}
\caption{Silueta y gap de transferencia por $k$ (media de 20 federaciones de FEMNIST).}
\label{tab:k}
\end{table}

\textbf{Brazos.} Tres semillas de datos, con la misma federación y sonda en todos los brazos: \texttt{p3} ($k=2$, NSGA-III con $\lambda=1$, 420 evaluaciones), \texttt{rs} (mismos clusters y objetivos, búsqueda aleatoria, 420 evaluaciones) y \texttt{base} ($C=10$, NSGA-II con (accuracy media, tamaño), 390 evaluaciones). Receta: 20 rondas $\times$ 2 épocas, Adam $\mathrm{lr}=10^{-3}$, evaluador retenido; cada modelo del hall of fame se reentrena una vez con los 100 nodos y se mide en test. La métrica primaria es el hipervolumen de los modelos reentrenados en (accuracy en test, tamaño normalizado).

\begin{table}[h]
\centering\small
\begin{tabular}{lccc}
\toprule
 & \texttt{p3} & \texttt{base} & \texttt{rs}\\
\midrule
hipervolumen & 0.515 & 0.518 & 0.499\\
mejor $\acc-\lambda\cdot$tamaño & 0.415 & 0.423 & 0.398\\
accuracy del mejor modelo & 0.514 & 0.516 & 0.498\\
tamaño medio del conjunto (normalizado) & 0.078 & 0.060 & 0.092\\
modelos en el conjunto final & 12 & 4--7 & 12\\
entrenamientos locales en la búsqueda & 16\,800 & 78\,000 & 16\,800\\
tiempo por corrida (h) & 0.98 & 1.28 & 0.95\\
\bottomrule
\end{tabular}
\caption{Resultado final (media de 3 semillas; test, modelos reentrenados con los 100 nodos).}
\label{tab:p3}
\end{table}

\begin{table}[h]
\centering\small
\begin{tabular}{llccc}
\toprule
contraste & métrica & media & IC95 & por semilla\\
\midrule
\texttt{p3}$-$\texttt{base} & hipervolumen & $-0.0024$ & $[-0.042,+0.019]$ & $+0.019,-0.042,+0.016$\\
\texttt{p3}$-$\texttt{base} & mejor $\acc-\lambda\cdot$tamaño & $-0.0072$ & $[-0.016,+0.0003]$ & $-0.016,+0.0003,-0.006$\\
\texttt{p3}$-$\texttt{base} & accuracy del mejor & $-0.0016$ & $[-0.051,+0.026]$ & $+0.021,-0.051,+0.026$\\
\texttt{p3}$-$\texttt{rs} & hipervolumen & $+0.0162$ & $[-0.032,+0.049]$ & $+0.049,+0.032,-0.032$\\
\texttt{p3}$-$\texttt{rs} & mejor $\acc-\lambda\cdot$tamaño & $+0.0170$ & $[-0.005,+0.050]$ & $+0.050,-0.005,+0.006$\\
\texttt{p3}$-$\texttt{rs} & accuracy del mejor & $+0.0158$ & $[-0.040,+0.050]$ & $+0.050,+0.038,-0.040$\\
\bottomrule
\end{tabular}
\caption{Contrastes pareados por semilla (diferencia media e IC95 bootstrap sobre 3 semillas, descriptivo).}
\label{tab:p3c}
\end{table}

Todos los intervalos incluyen 0, así que, por la regla fijada de antemano, no se declara ventaja de ningún brazo; ni siquiera en $\acc-\lambda\cdot$tamaño, el criterio que optimiza \texttt{p3}, supera a \texttt{base}. Lo único claro es el coste: \texttt{p3} usa 4.6 veces menos entrenamientos locales (2 representantes por evaluación frente a 10) con resultados que no se distinguen, con solo 3 semillas.

\textbf{¿Hubo evolución?} Sí en la población y poco en el mejor individuo (Tabla~\ref{tab:evo}; solo hijos desde la segunda generación, con $S=$ accuracy estimada $-$ tamaño). En \texttt{p3} la calidad media de los individuos nuevos sube $+0.10$ de la primera a la última generación (positiva en las 3 semillas) y en \texttt{rs} $+0.004$; la diferencia \texttt{p3}$-$\texttt{rs}, con los mismos nodos representantes en ambos brazos en cada generación, pasa de $\approx$0.02 (generación 0) a $\approx$0.12 (generación 3). \texttt{base} mejora $+0.18$ (3 de 3), casi todo en las tres primeras generaciones, y su tamaño medio baja de 0.12 a 0.06. Sin embargo, el mejor $S$ acumulado sube parecido en \texttt{p3} (0.314 a 0.414) y en \texttt{rs} (0.28 a 0.37): parte de la subida de \texttt{rs} es solo tener más muestras aleatorias. Con arquitecturas aleatorias malas en promedio (accuracy estimada de 0.14 en la generación 0), unos cientos de muestras al azar ya incluyen alguna muy buena: la evolución aumenta la proporción de buenas, pero ayuda poco a encontrar la mejor.

\begin{table}[h]
\centering\small
\begin{tabular}{lccc}
\toprule
generación & \texttt{p3} & \texttt{rs} & \texttt{base}\\
\midrule
0 & 0.025 & 0.007 & 0.105\\
1 & 0.040 & $-0.006$ & --\\
2 & 0.146 & 0.051 & --\\
3 & 0.127 & 0.010 & 0.306\\
última & 0.127 & 0.010 & 0.282 (gen.\ 19)\\
\bottomrule
\end{tabular}
\caption{$S$ medio de los individuos nuevos por generación (media de 3 semillas).}
\label{tab:evo}
\end{table}

\section{Límites}
\begin{itemize}[leftmargin=*,itemsep=2pt]
\item Los gaps de transferencia son de accuracy absoluta con modelos débiles (nodos de 90 a 350 muestras); en FEMNIST el efecto es de pocos puntos.
\item Una sola arquitectura de medición en los estudios de transferencia, un solo \texttt{min\_samples} en DBSCAN; solo CIFAR-10 y FEMNIST.
\item Con Dirichlet la similitud funcional depende casi solo de las etiquetas. En FEMNIST el coeficiente de $\approx$1 punto sobre Hellinger es exploratorio y no identifica qué captura $\Delta$.
\item Arquitectura $\times$ nodo: 24 arquitecturas aleatorias del espacio v2, una receta y solo FEMNIST; test por nodo muy pequeño (36 muestras); el IC bootstrap remuestrea nodos, no arquitecturas; no se identificó qué arquitecturas o nodos generan la interacción.
\item P3: 3 semillas (intervalos de $\pm0.03$--$0.05$), un $\lambda$, 4 generaciones en \texttt{p3} y \texttt{rs}; evaluaciones casi iguales (420 y 390) pero entrenamientos locales distintos; los conjuntos finales tienen tamaños distintos (el hall of fame de \texttt{base} es su frente no dominado, de 4 a 7 modelos), lo que complica comparar hipervolúmenes; la mejora por generación se mide sobre la accuracy estimada (pocos nodos, ruidosa), no sobre test; la penalización común elimina el frente accuracy-tamaño. Como la silueta eligió $k=2$, no se probó el caso de muchos objetivos.
\end{itemize}

\end{document}
